\documentclass[10pt]{article}
\usepackage[margin=20mm]{geometry}
\usepackage{amsmath,amssymb,mathtools,fontspec,unicode-math,xurl,hyperref}
\IfFontExistsTF{DejaVu Serif}{\setmainfont{DejaVu Serif}}{\setmainfont{Latin Modern Roman}}
\setmathfont{latinmodern-math.otf}
\setlength{\parindent}{0pt}\setlength{\parskip}{6pt}\sloppy
\title{Initial technical review of
Spectral analysis and \(k\)-spine decomposition of inhomogeneous branching Brownian motions. Genealogies in fully pushed fronts}
\author{AI-assisted technical review}
\date{}
\begin{document}
\maketitle
Abstract: The acquired journal version contains a false boundary-uniform moment estimate in the central survival-probability argument. A starting point at fixed distance from the moving upper boundary has reproductive weight of order \(N\); its initial two-spine contribution over time \(L^2\) diverges as \(N/L^2\). An associated lower-bound statement also has an impossible uniformity in the time interval. This is a Major defect in the stated derivation, without establishing that the fixed-starting-point main theorem is false. Three further, bounded findings concern an extension to the semipushed regime, a missing factor of two in the Green kernel, and an incorrect high-frequency compactness argument. Exact repairs are supplied for the kernel normalization and the high-frequency norm estimate. The report does not claim a complete replacement proof of the main results.\par
\section*{Reviewed object, assumptions, and positive checks}
This is the initial technical review of M25-0514, by Emmanuel Schertzer and Julie Tourniaire, The Annals of Probability 53(4) (2025), 1382–1433, DOI 10.1214/24-AOP1736 (https://doi.org/10.1214/24-AOP1736) [ST]. The acquired version of record has 52 physical pages, without an additional cover. Printed page \(p\) is physical page \(p-1381\). Its SHA256 is\par
> 00000TOKEN.\par
Scientific acceptance remains pending independent review.\par
Equation reference: eq:weights. The model is dyadic branching Brownian motion with unit Brownian variance, drift \(-\mu\), killing at zero, and rate \(r(x)=(1+W(x))/2\). Assumptions (A1)–(A2), page 1383, require a nonnegative \(C^1\) function \(W\) supported in \([0,1]\). For the process also killed at \(L\), let \(v_L\) be the positive Dirichlet principal eigenfunction of \[\tfrac12 v_L''+\tfrac12 Wv_L=\lambda_1(L)v_L,\qquad v_L(1)=1.\] Write \[\lambda_\infty=\lim_{L\to\infty}\lambda_1(L),\quad
 \beta=\sqrt{2\lambda_\infty},\quad
 \mu=\sqrt{1+2\lambda_\infty},\quad
 \alpha=\frac{\mu+\beta}{\mu-\beta},\quad
 w_L=\lambda_\infty-\lambda_1(L).\] The pushed assumption \(H_p\) is \(\lambda_\infty>0\). The main theorems assume \(H_{fp}\), namely \(\lambda_\infty>1/16\), equivalently \(\mu<3\beta\) or \(\alpha>2\). The cutoff in (19) and the weights in (20), page 1393, are \[
 L=L_N=\frac{\log N}{\mu-\beta},\qquad
 H_L(x):=h(0,x)=\frac{e^{\mu x}v_L(x)}
                    {\widetilde c\|v_L\|_2^2},\qquad
 h(s,x)=e^{w_Ls}H_L(x),\] where \(\widetilde c^{-1}=\int_0^\infty e^{-\mu x}v_\infty(x)\,dx\). Here \(v_\infty(x)=e^{-\beta(x-1)}\) for \(x\geq1\), and \(I:=\|v_\infty\|_2^2\in(0,\infty)\).\par
Several structural calculations check correctly. For \(A=\frac12\partial_{xx}-\mu\partial_x+r\), direct substitution gives \(AH_L=-w_LH_L\). Consequently \[(\partial_s+A)h=0,\qquad
 H_L^{-1}A(H_L f)+w_Lf
   =\tfrac12f''+\frac{v_L'}{v_L}f'.\] Thus the conservative spine has invariant density \(\Pi_L=v_L^2/\|v_L\|_2^2\); its probability current vanishes since \(\frac12\Pi_L'=(v_L'/v_L)\Pi_L\). Also \(r h_\infty\Pi_\infty\) has tail proportional to \(e^{(\mu-3\beta)x}\). Its integrability holds precisely when \(\mu<3\beta\), giving the stated threshold for the finite variance constant \(\Sigma^2=2\int r h_\infty\Pi_\infty\). Near zero the integrand is of order \(x^3\). These checks support the basic transform and the fully pushed integrability criterion; they do not supply the missing uniform estimate identified next.\par
\section*{A boundary-uniform estimate fails in the survival proof}
Problem 1 (Major: the initial boundary contribution is not uniformly bounded).  Lemma 15, page 1424 (physical 43), claims a positive constant \(C_1=C_1(\lambda_\infty)\), independent of \(0<T_1<T_2\), such that \[u(tN,x)\geq \frac{C_1H_L(x)}{1+tN},
 \qquad x\in[0,L],\quad t\in[T_1,T_2],\] where \(u(s,x)=\mathbb P_x(Z_s^L>0)\). Its proof on page 1425 asserts an interval-independent bound on the two-spine expectation. More decisively, the proof of Corollary 8, page 1427 (physical 46), after (94), claims that \[Q_x^{2,L^2}\big[r(\zeta_v)h(|v|,\zeta_v)\big]\] is uniformly bounded in \(N\) and \(x\). Both uniformity claims fail under the actual fully pushed hypotheses.\par
\subsection*{An exact admissible boundary sequence}
Equation reference: eq:gap, eq:boundary. Fix any potential satisfying (A1)–(A2) and \(H_{fp}\), and fix \(d>0\). For large \(N\), \(x_N=L_N-d>1\). The explicit tail (43) gives \[v_L(L-d)=\frac{\sinh(\beta_Ld)}
                   {\sinh(\beta_L(L-1))},
 \qquad \beta_L=\sqrt{2\lambda_1(L)}.\] The valid eigenvalue-gap calculation (45) implies \((\beta-\beta_L)L\to0\). One can independently obtain its leading term from the endpoint derivative formula \[\lambda_1'(L)=\frac{v_L'(L)^2}{2\|v_L\|_2^2}
 \sim\frac{2\beta^2e^{2\beta}}{I}e^{-2\beta L},\] then integrate from \(L\) to infinity. Hence \[
 w_L\sim \frac{\beta e^{2\beta}}{I}e^{-2\beta L}
       =\frac{\beta e^{2\beta}}{I}N^{-(\alpha-1)}.\] Together with \(\|v_L\|_2^2\to I\), the explicit tail yields \[
 \frac{H_L(L-d)}{N}
 \longrightarrow
 K_d:=\frac{2e^\beta}{\widetilde c I}
          e^{-\mu d}\sinh(\beta d)>0 .\] This is a limit of the precise normalized weight in (20); there is no freedom to remove the factor \(N\).\par
The assumptions are nonempty. For example \(W(x)=200(1-x)^2\) on \([0,1]\), zero elsewhere, is nonnegative and \(C^1\). Testing the Rayleigh quotient with \(\sin(2\pi x)\) on \([0,1/2]\), extended by zero, gives \(\lambda_\infty\geq25-2\pi^2>1/16\). The argument above works for every admissible fully pushed potential, rather than relying on an approximate eigenvalue for this example.\par
Contradiction to Lemma 15 as quantified. If the positive \(C_1\) is independent of the time interval, choose a fixed \(t<C_1K_d\) and then choose \(0<T_1<t<T_2\). By [eq:boundary] (\#eq:boundary), \(C_1H_L(x_N)/(1+tN)\to C_1K_d/t>1\). The claimed probability lower bound therefore exceeds one for large \(N\), which is impossible. Allowing \(C_1\) to depend on \(T_1\) removes this particular contradiction, but changes the statement and does not by itself justify the subsequent limit as the initial macroscopic time tends to zero in (81). ◻\par
\subsection*{Direct contradiction to the short-time two-spine bound}
Equation reference: eq:q2. The tree construction on pages 1394–1395 and Corollary 1, page 1404, make the sole branch time of a two-spine tree uniform on \([0,T]\). Conditional on that forward time \(s\), the branching location has the one-spine law at \(s\). Therefore, exactly, \[
 Q_x^{2,T}[r(\zeta_v)h(|v|,\zeta_v)]
 =\frac1T\int_0^T
       \mathbb E_x^Q[r(\zeta_s)h(s,\zeta_s)]\,ds .\] This is the unaccelerated \(Q^{2,T}\) used in (94); the accelerated notation elsewhere does not change this identity.\par
For \(y>1\), the spine drift is \[\frac{v_L'(y)}{v_L(y)}
   =-\beta_L\coth\big(\beta_L(L-y)\big).\] It is uniformly bounded in absolute value on \([L-2d,L-d/2]\) for all large \(N\). Starting from \(x_N=L-d\), there are \(s_0>0\) and \(p>0\), independent of \(N\), such that the spine stays in that interval throughout \([0,s_0]\) with probability at least \(p\). For completeness, stop the SDE at the interval endpoints, bound its drift by \(B\), choose \(Bs_0<d/8\), and use the Brownian event \(\sup_{s\leq s_0}|B_s|<d/8\), which has strictly positive probability. On this event the stopped path cannot reach an endpoint.\par
Equation reference: eq:diverge. Uniformly over that interval, the same hyperbolic calculation as [eq:boundary] (\#eq:boundary) gives \(H_L(y)\geq cN\) for some \(c>0\). Also \(r(y)=1/2\) there and \(h(s,y)\geq H_L(y)\). Applying [eq:q2] (\#eq:q2) with \(T=L^2\) proves \[
 Q_{x_N}^{2,L^2}[r(\zeta_v)h(|v|,\zeta_v)]
 \geq \frac{pcs_0}{2}\frac{N}{L^2}
 \longrightarrow\infty .\] Thus the actual estimate invoked after (94) is false, independently of the interpretation of the constants in Lemma 15. For \(T=tN\), the same initial contribution is at least \(pcs_0/(2t)\); this also prevents a single bound valid for all arbitrarily small positive macroscopic times.\par
\subsection*{Dependency, alternative readings, and repair boundary}
The disputed estimate bounds the Bonferroni pair term in (94). It is used to obtain (96), the lower half of Corollary 8’s all-\(x\) profile approximation. Corollary 8 then enters Lemma 17 and the survival differential equation; page 1428 uses these results to prove Theorem 1. Theorem 1 in turn normalizes the conditional moments and the genealogy limits in Section 7.\par
The counterexample uses interior starting points at distance \(d\) from the boundary, not a particle started at the absorbing endpoint. It also concerns a positive expectation; there is no cancellation that can remove the initial contribution. Mixing on a scale much shorter than \(N\) controls later times but does not erase the positive integral over \([0,s_0]\). Lemma 13 itself allows a factor \(t^{-(k-1)}\); substituting \(t=L^2/N\) does not give the claimed constant bound.\par
This is a demonstrated invalid step in a core derivation and merits Major severity. It is not a counterexample to Theorem 1 for fixed \(x>0\), nor a proof that the final Yaglom or genealogy limits are false. A possible repair must control the initial contribution with its spatial dependence, restrict an intermediate uniformity range, or justify a weighted integration that makes the boundary region negligible. The report does not claim to have completed that replacement argument. It also does not infer the falsity of Corollary 8 solely from the failure of its proof.\par
\section*{The stated pushed range is too broad for the time normalization}
Problem 2 (Minor: equation (46) and the second limit of Corollary 2 fail under \(H_p\)).  Lemma 5, page 1407 (physical 26), assumes only \(H_p\) and claims, for \(0<t\lesssim N\), \[h(t,x)=\bigl(1+O(e^{(\mu-3\beta)L})\bigr)h(0,x).\] Corollary 2, pages 1408–1409 (physical 27–28), likewise assumes only \(H_p\), and asserts convergence for bounded continuous tests of \((x,h(tN,x))\), uniformly over bounded macroscopic time intervals. These statements fail in the semipushed range.\par
Choose an admissible \(W\) with \(\lambda_\infty=1/32\). Such a potential exists in the family \(W_a=a(1-x)^2\mathbf1_{[0,1]}(x)\): the supremum of the Rayleigh quotient is nonnegative, starts at zero, is Lipschitz in \(a\) with constant at most \(1/2\), and tends to infinity by any fixed nonzero test function in \((0,1)\). Continuity gives the asserted value. Then \[\beta=\tfrac14,\quad \mu=\tfrac{\sqrt{17}}4,\quad
 \alpha=\frac{\sqrt{17}+1}{\sqrt{17}-1}\in(1,2).\] At \(t=N\), [eq:gap] (\#eq:gap) gives \[w_LN\sim C N^{2-\alpha}\longrightarrow\infty,\qquad C>0.\] But the exact time factor in [eq:weights] (\#eq:weights) is \(e^{w_LN}\). It cannot equal \(1+O(N^{2-\alpha})\), whereas \(e^{(\mu-3\beta)L}=N^{2-\alpha}\). The error is an invalid linearization of an exponential whose argument diverges.\par
There is also a bounded continuous test contradicting the second limit in Corollary 2. Take \(k=1\), \(t_1=1\), and \(f(x,z)=\chi(x)e^{-z}\), where \(\chi\) is a nonnegative, nonzero continuous function with compact support in \((0,\infty)\). On its support \(H_L\to h_\infty>0\) uniformly, so \[\int \chi(x)e^{-h(N,x)}\Pi_L(dx)\longrightarrow0.\] The asserted limit is instead \(\int\chi(x)e^{-h_\infty(x)}\Pi_\infty(dx)>0\). All tests and parameter restrictions in the claimed corollary are satisfied.\par
The first, time-independent convergence of \(\Pi_L\) is not refuted. Nor is this an example in the fully pushed main-theorem regime. The introductory sentence of Section 4, page 1405, expressly says the results of Sections 4.1–4.2 do not require \(H_{fp}\), so the stronger assumption cannot silently be inserted as the original statement. Requiring \(H_{fp}\) for the \(N\)-time normalization, or restricting time so that \(w_Lt=o(1)\), is a bounded repair. Keeping \(e^{w_Lt}\) exactly is another option when formulating a different limit. The minor severity reflects this range restriction and the main results’ explicit \(H_{fp}\) hypothesis.\par
\section*{The Green kernel and Wronskian have a factor-of-two discrepancy}
Problem 3 (Minor: incorrect normalization in (63) and Lemma 10).  Section 4.3, page 1413 (physical 32), defines \[\omega_\lambda=
 \psi_\lambda(1)\varphi_\lambda'(1)
   -\psi_\lambda'(1)\varphi_\lambda(1),\] where \(\varphi_\lambda,\psi_\lambda\) solve \(\frac12u''+\frac12Wu=\lambda u\), with \(\varphi_\lambda(0)=0\), \(\psi_\lambda(L)=0\), and the normalizations (62). Proposition 13, page 1414, gives the resolvent kernel as \(\omega_\lambda^{-1}\) times the product of fundamental solutions. For the stated diffusion coefficient \(1/2\), that coefficient must be \(2\omega_\lambda^{-1}\). Lemma 10, pages 1416–1417, inherits the same factor discrepancy.\par
Equation reference: eq:green. Exact jump and Wronskian checks. Put \(\lambda=\lambda_\infty+\xi\), \(\xi>0\), and remove the drift factor by writing \[R_\lambda(x,y)=e^{\mu(y-x)}G_\xi(x,y).\] It must solve \[\left(\lambda-\tfrac12\partial_{yy}-\tfrac12W(y)\right)
 R_\lambda(x,y)=\delta_x(y).\] Integrating across \(y=x\) shows that its derivative jump must equal \(-2\). The continuous product kernel \(\varphi_\lambda(\min(x,y))\psi_\lambda(\max(x,y))/\omega_\lambda\) has derivative jump \(-1\), by the displayed definition of \(\omega_\lambda\). Therefore the exact correction is \[
 G_\xi(x,y)=\frac{2e^{\mu(x-y)}}{\omega_{\lambda_\infty+\xi}}
 \varphi_{\lambda_\infty+\xi}(\min(x,y))
 \psi_{\lambda_\infty+\xi}(\max(x,y)).\] Multiplying fundamental solutions by arbitrary constants leaves this conclusion unchanged.\par
Equation reference: eq:wronskian. There is an independent exact identity for the Wronskian. Let \(v=v_L\) and \(F(y)=\psi_\lambda(y)v'(y)-\psi_\lambda'(y)v(y)\). Then \[F'(y)=2(\lambda_1-\lambda)\psi_\lambda(y)v(y),\quad
 F(L)=0,\quad F(0)=\omega_\lambda,\] where the last equality uses \(\varphi_\lambda'(0)=v_L'(0)\). Consequently \[
 \omega_{\lambda_\infty+\xi}
   =2(\xi+w_L)\int_0^L\psi_{\lambda_\infty+\xi}(y)v_L(y)\,dy .\] Under \(H_{fp}\) and \(N^{-1}\lesssim\xi\ll L^{-1}\), the elementary continuous-dependence estimates of Lemma 9 give \[\omega_{\lambda_\infty+\xi}
 =2(\xi+w_L)\bigl(\|v_L\|_2^2+O(\xi L)\bigr)
 \sim 2\xi I .\] Lemma 10 instead states leading term \(\xi I\). Retaining \(\xi+w_L\) also avoids suppressing a relative error \(w_L/\xi\) that need not be \(O(\xi L)\). For example, \(\xi=N^{-1}\) and \(2<\alpha<3\) give \(w_L/\xi\asymp N^{-(\alpha-2)}\), larger than \(L/N\). ◻\par
The identity defining \(G_\xi\) is the time integral of the Lebesgue-density particle kernel, and the governing PDE has diffusion coefficient \(1/2\). Thus neither a speed-measure convention nor a rescaling of the Wronskian can reconcile (63) with those definitions. As a separate algebraic calibration, for \(W=0\), \(\varphi(y)=\sinh(ky)/k\) and \(\psi(y)=\sinh(k(L-y))\), the printed jump is exactly \(-1\). That calibration alone is not asserted to satisfy \(H_p\); the general jump argument above applies to the actual admissible potentials.\par
The paper uses this resolvent in Lemmas 11–12 to obtain order bounds on escape and moment integrals. The repaired Wronskian still has \(\omega^{-1}=O(N)\) at \(\xi=(TN)^{-1}\) for fixed \(T>0\), so this normalization error alone does not refute those order estimates or the final \(\Sigma^2\). Equations [eq:green] (\#eq:green)–[eq:wronskian] (\#eq:wronskian) give a direct local repair. The cited handbook page was not accessible in the bounded reference search; the proof here uses the paper’s own PDE and does not attribute an error to that external source.\par
\section*{A false positive-minimum argument has a local replacement}
Problem 4 (Minor: the compactness argument in Lemma 8 fails at zero).  Lemma 8, pages 1410–1411 (physical 29–30), seeks \(\|u_k\|_2^2\geq c_8(1\wedge|\lambda_k|^{-1})\), where \(u_\lambda(0)=0\), \(u_\lambda'(0)=1\). After (55), it rescales the interval by \(1/|\lambda|\) and asserts that the function \(g_2\) on \([-1,0]\) attains a strictly positive minimum. Its continuous extension at zero is actually zero. The asserted norm estimate can nevertheless be proved by a different argument.\par
Equation reference: eq:g2. Write \(\lambda=-a\), \(a\to\infty\), and \(z_a(t)=a u_{-a}(t/a)\). Since \(u_{-a}''=(-2a-W)u_{-a}\), \[z_a''(t)=-
 \left(\frac2a+\frac{W(t/a)}{a^2}\right)z_a(t),\qquad
 z_a(0)=0,\quad z_a'(0)=1.\] For bounded \(W\), elementary continuous dependence of this integral equation gives \(z_a\to t\) uniformly on \([0,1]\). The integrand used to define \(g_2(-1/a)\) is exactly \(u_{-a}(t/a)^2\), because \(u_\lambda=\rho_\lambda\sin\theta_\lambda\). Therefore \[
 a^2g_2(-1/a)=\int_0^1z_a(t)^2\,dt\longrightarrow\frac13.\] There is no positive minimum. Moreover the lower integral in (55) is only asymptotic to \(1/(3a^3)\), so that particular lower bound cannot yield the desired \(1/a\) scale.\par
For an actual admissible potential, the independent calculation uses \(W=200(1-x)^2\) on \([0,1]\) as above. An 80-digit power-series solution of the scaled equation gives the following values; degree 40 and degree 60 evaluations agree to the shown precision. They corroborate [eq:g2] (\#eq:g2), rather than replacing its proof. \[\begin{array}{c|cccc}
 a&100&1000&10000&10^6\\ \hline
 a^2g_2(-1/a)&0.3306877972&0.3331867085&
              0.3333198669&0.3333332000
\end{array}\] No certified floating-point interval bound is claimed.\par
A replacement proof of the norm estimate. For \(\lambda=-s^2/2\), variation of constants gives \[u_\lambda(x)=\frac{\sin(sx)}s
   -\frac1s\int_0^x\sin(s(x-y))W(y)u_\lambda(y)\,dy .\] Let \(M=\|W\|_\infty\). Gronwall’s inequality on \([0,1]\) gives \[\|u_\lambda\|_\infty\leq \frac{e^{M/s}}s,\qquad
 \left\|u_\lambda-\frac{\sin(s\,\cdot)}s\right\|_\infty
       \leq\frac{Me^{M/s}}{s^2}.\] For sufficiently large \(s\), \[\left\|\frac{\sin(s\,\cdot)}s\right\|_{L^2(0,1)}
 =\frac1s\left(\frac12-\frac{\sin(2s)}{4s}\right)^{1/2}
 \geq\frac1{2s},\] and the error norm is at most \(1/(4s)\). Thus \[\|u_\lambda\|_{L^2(0,1)}^2
    \geq\frac1{16s^2}=\frac1{32|\lambda|}.\] On the remaining compact range of \(\lambda\leq\lambda_\infty\), continuity of the ODE solution and the nonzero initial derivative give a positive minimum of \(\int_0^1u_\lambda^2\). Since \(L\geq1\), the full norm over \([0,L]\) is at least this norm. Combining the two ranges proves Lemma 8’s claimed order, with a constant depending on the fixed potential. ◻\par
Lemma 8 supports the normalized-eigenfunction bounds in Proposition 10 and hence the heat-kernel estimate in Proposition 12. The replacement above restores the needed order without the false compactness assertion. Accordingly this is a proof defect with a complete local repair, not a counterexample to the norm lemma or by itself to the main theorem.\par
\section*{Coverage, counterevidence, and reviewed-version limits}
The primary reviewer read the extracted scientific text of all 52 physical pages. The review followed the model assumptions, main theorem range, the many-to-one and many-to-few weights, spectral normalization, moment estimates, survival argument, and their use in conditional metric-space convergence. The four findings above include all hypotheses used by their counterarguments. Reading and checking these dependency paths does not amount to independent certification of every cited theorem or every statement in the paper.\par
All source pages bearing the reported evidence were also inspected as native images: physical pages 1–6, 12–15, 17–18, 20–24, 26–30, 32–33, 35–36, 39–41, 43–47, and 49–52. The explicit hyperbolic tail, branch-time convention, time quantifiers, Green-kernel coefficient, and high-frequency minimum assertion were checked against those page images. Other extracted pages were read as text, not claimed as image inspections. Mechanical extraction and rendering are recorded separately from actual viewing.\par
Bounded searches on 29 October 2026 checked the authors’ publication listings, the arXiv record, and title/DOI-specific correction and supplement queries. The public arXiv record lists version 3 dated 18 October 2024 [arxiv]; the author listing identifies the published article [authorpage]. No target-specific correction was found in those checks. This is not a claim that no correction exists. The publisher landing-page route encountered an access/security limit. The acquired version of record remains the reviewed object; the arXiv file was not silently substituted.\par
The cited Borodin–Salminen DOI resolves to the first edition, copyright 1996, with a 2012 electronic publication date [BSmetadata]. Its public chapter information did not provide the cited page 19. That source page was neither read nor viewed. No claim about the handbook’s own normalization is made. No author code was executed; no lack of author code is treated as a mathematical defect.\par
Problem 1 remains a substantive gap in the printed core derivation. Problems 2–4 have explicitly bounded impacts, with local repairs where established. These findings do not constitute a proof of failure of the fully pushed fixed-\(x\) limit theorems, a complete revised proof, or an independent acceptance decision.\par
9 E. Schertzer and J. Tourniaire, Spectral analysis and \(k\)-spine decomposition of inhomogeneous branching Brownian motions. Genealogies in fully pushed fronts, Ann. Probab. 53 (2025), 1382–1433. doi:10.1214/24-AOP1736 (https://doi.org/10.1214/24-AOP1736). Acquired version of record, hash specified in Section 1. E. Schertzer and J. Tourniaire, arXiv metadata for Spectral analysis and \(k\)-spine decomposition of inhomogeneous branching Brownian motions. Genealogies in fully pushed fronts, https://arxiv.org/abs/2301.01697, version 3 dated 18 October 2024. Metadata checked 29 October 2026. J. Tourniaire, publication listing, https://jtourniaire.pages.math.cnrs.fr/. Checked 29 October 2026 for publication and correction context. A. N. Borodin and P. Salminen, Handbook of Brownian Motion—Facts and Formulae, publisher metadata for first edition, copyright 1996, electronic publication 6 December 2012, https://doi.org/10.1007/978-3-0348-7652-0. Metadata only; cited page 19 was unavailable.\par
\end{document}
